% GENERATED FILE. Edit canonical lessons, then run npm run generate:books.
% canonical-source-hash: fnv1a64:cb04675f7c388964
% canonical-lessons: UR-C356-islah, UR-C356-meyar, UR-C356-saliqa, UR-C356-sudharna, UR-C356-nikharna

\chapter{Making Things Better}
\label{ch:ur-making-things-better}
\hlchaptermodality{356}

\begin{chapteropening}
\textbf{By the end of this chapter:} \emph{I can say how a place or a person improved.}

Complete the last lesson of chapter 356: I can say how a place or a person improved.
\end{chapteropening}

\section[\texorpdfstring{\ur{اصلاح} (islāh)}{اصلاح (islāh)}]{\ur{اصلاح} (islāh) — a reform; a correction}

\label{lesson:UR-C356-islah}



\begin{quote}
Before the new one: say the Urdu for to commend, then the Urdu for to endure.
\end{quote}

\subsection*{You'll want to know: \ur{اصلاح}}
\textbf{\ur{اصلاح}} — \emph{islāh} — \textquotedblleft{}a reform; a correction\textquotedblright{}.

\begin{quote}
\ur{اسکولوں} \ur{میں} \ur{اصلاح} \ur{کی} \ur{ضرورت} \ur{ہے۔}
\end{quote}

\begin{quote}
\emph{iskūloṅ meṅ islāh kī zarūrat hai.}
\end{quote}

\begin{quote}
The schools need reform.
\end{quote}

Feminine.

\begin{cousinweb}
\textbf{\ur{اصلاح}} (\emph{islāh}) is usually traced, through Persian, to Arabic \emph{iṣlāḥ}, \textquotedblleft{}putting right\textquotedblright{}.
\end{cousinweb}

\begin{grammarlens}[title={What you've built}]
One more word for giving reasons and opinions.
\end{grammarlens}

\subsection*{Your turn}
\begin{itemize}
\raggedright
  \item \emph{Say it:} \emph{islāh}
  \item \emph{Say it:} \emph{islāh}, once more
  \item \emph{Say it:} \emph{iskūloṅ meṅ islāh kī zarūrat hai.} — The schools need reform.
  \item \emph{Recall:} say the Urdu for to commend, then the Urdu for to endure, then say \emph{islāh} again
\end{itemize}

\begin{tcolorbox}[breakable,colback=teal!4,colframe=teal!35!black,title={Before you move on}]
What does \ur{اصلاح} mean? (\textquotedblleft{}A reform; a correction\textquotedblright{}.) Say it once more.
\end{tcolorbox}
\section[\texorpdfstring{\ur{معیار} (meyār)}{معیار (meyār)}]{\ur{معیار} (meyār) — a standard (of quality)}

\label{lesson:UR-C356-meyar}



\begin{quote}
Before the new one: say the Urdu for to endure, then the Urdu for a reform.
\end{quote}

\subsection*{You'll want to know: \ur{معیار}}
\textbf{\ur{معیار}} — \emph{meyār} — \textquotedblleft{}a standard (of quality)\textquotedblright{}.

\begin{quote}
\ur{اس} \ur{دکان} \ur{کا} \ur{معیار} \ur{اچھا} \ur{ہے۔}
\end{quote}

\begin{quote}
\emph{is dukān kā meyār achchhā hai.}
\end{quote}

\begin{quote}
This shop keeps a good standard.
\end{quote}

Masculine.

\begin{cousinweb}
\textbf{\ur{معیار}} (\emph{meyār}) is usually traced, through Persian, to Arabic \emph{miʿyār}, \textquotedblleft{}a measure, a touchstone\textquotedblright{}.
\end{cousinweb}

\begin{grammarlens}[title={What you've built}]
One more word for giving reasons and opinions.
\end{grammarlens}

\subsection*{Your turn}
\begin{itemize}
\raggedright
  \item \emph{Say it:} \emph{meyār}
  \item \emph{Say it:} \emph{meyār}, once more
  \item \emph{Say it:} \emph{is dukān kā meyār achchhā hai.} — This shop keeps a good standard.
  \item \emph{Recall:} say the Urdu for to endure, then the Urdu for a reform, then say \emph{meyār} again
\end{itemize}

\begin{tcolorbox}[breakable,colback=teal!4,colframe=teal!35!black,title={Before you move on}]
What does \ur{معیار} mean? (\textquotedblleft{}A standard (of quality)\textquotedblright{}.) Say it once more.
\end{tcolorbox}
\section[\texorpdfstring{\ur{سلیقہ} (salīqā)}{سلیقہ (salīqā)}]{\ur{سلیقہ} (salīqā) — a knack for doing things well; good taste}

\label{lesson:UR-C356-saliqa}



\begin{quote}
Before the new one: say the Urdu for a reform, then the Urdu for a standard.
\end{quote}

\subsection*{You'll want to know: \ur{سلیقہ}}
\textbf{\ur{سلیقہ}} — \emph{salīqā} — \textquotedblleft{}a knack for doing things well; good taste\textquotedblright{}.

\begin{quote}
\ur{اسے} \ur{کھانا} \ur{پکانے} \ur{کا} \ur{سلیقہ} \ur{ہے۔}
\end{quote}

\begin{quote}
\emph{use khānā pakāne kā salīqā hai.}
\end{quote}

\begin{quote}
She has a knack for cooking.
\end{quote}

Masculine.

\begin{cousinweb}
\textbf{\ur{سلیقہ}} (\emph{salīqā}) is usually traced, through Persian, to Arabic \emph{salīqa}, \textquotedblleft{}a natural gift, a disposition\textquotedblright{}.
\end{cousinweb}

\begin{grammarlens}[title={What you've built}]
One more word for giving reasons and opinions.
\end{grammarlens}

\subsection*{Your turn}
\begin{itemize}
\raggedright
  \item \emph{Say it:} \emph{salīqā}
  \item \emph{Say it:} \emph{salīqā}, once more
  \item \emph{Say it:} \emph{use khānā pakāne kā salīqā hai.} — She has a knack for cooking.
  \item \emph{Recall:} say the Urdu for a reform, then the Urdu for a standard, then say \emph{salīqā} again
\end{itemize}

\begin{tcolorbox}[breakable,colback=teal!4,colframe=teal!35!black,title={Before you move on}]
What does \ur{سلیقہ} mean? (\textquotedblleft{}A knack for doing things well; good taste\textquotedblright{}.) Say it once more.
\end{tcolorbox}
\section[\texorpdfstring{\ur{سدھارنا} (sudhārnā)}{سدھارنا (sudhārnā)}]{\ur{سدھارنا} (sudhārnā) — to improve; to put right}

\label{lesson:UR-C356-sudharna}



\begin{quote}
Before the new one: say the Urdu for a standard, then the Urdu for a knack.
\end{quote}

\subsection*{You'll want to know: \ur{سدھارنا}}
\textbf{\ur{سدھارنا}} — \emph{sudhārnā} — \textquotedblleft{}to improve; to put right\textquotedblright{}.

\begin{quote}
\ur{اس} \ur{نے} \ur{اپنا} \ur{کام} \ur{سدھارا۔}
\end{quote}

\begin{quote}
\emph{us ne apnā kām sudhārā.}
\end{quote}

\begin{quote}
She improved her work.
\end{quote}

Past, with \ur{نے} (he or she …): \ur{اس} \ur{نے} \ur{گھر} \ur{سدھارا}, \emph{us ne ghar sudhārā}.

\begin{cousinweb}
\textbf{\ur{سدھارنا}} (\emph{sudhārnā}) is inherited: CDIAL lists the Hindi cousin \emph{sudhārnā}, \textquotedblleft{}to mend\textquotedblright{}, under \emph{*śuddhakāra}, \textquotedblleft{}making clean\textquotedblright{}, a form it reconstructs (CDIAL 12521).
\end{cousinweb}

\begin{grammarlens}[title={What you've built}]
One more word for giving reasons and opinions.
\end{grammarlens}

\subsection*{Your turn}
\begin{itemize}
\raggedright
  \item \emph{Say it:} \emph{sudhārnā}
  \item \emph{Say it:} \emph{sudhārnā}, once more
  \item \emph{Say it:} \emph{us ne apnā kām sudhārā.} — She improved her work.
  \item \emph{Recall:} say the Urdu for a standard, then the Urdu for a knack, then say \emph{sudhārnā} again
\end{itemize}

\begin{tcolorbox}[breakable,colback=teal!4,colframe=teal!35!black,title={Before you move on}]
What does \ur{سدھارنا} mean? (\textquotedblleft{}To improve; to put right\textquotedblright{}.) Say it once more.
\end{tcolorbox}
\section[\texorpdfstring{\ur{نکھارنا} (nikhārnā)}{نکھارنا (nikhārnā)}]{\ur{نکھارنا} (nikhārnā) — to polish a skill; to bring out the best in}

\label{lesson:UR-C356-nikharna}



\begin{quote}
Before the new one: say the Urdu for a knack, then the Urdu for to put right.
\end{quote}

\subsection*{You'll want to know: \ur{نکھارنا}}
\textbf{\ur{نکھارنا}} — \emph{nikhārnā} — \textquotedblleft{}to polish a skill; to bring out the best in\textquotedblright{}.

\begin{quote}
\ur{استاد} \ur{نے} \ur{اس} \ur{کا} \ur{ہنر} \ur{نکھارا۔}
\end{quote}

\begin{quote}
\emph{ustād ne us kā hunar nikhārā.}
\end{quote}

\begin{quote}
The teacher brought out the best in her skill.
\end{quote}

Past, with \ur{نے} (he or she …): \ur{اس} \ur{نے} \ur{اپنا} \ur{ہنر} \ur{نکھارا}, \emph{us ne apnā hunar nikhārā}.

\begin{cousinweb}
\textbf{\ur{نکھارنا}} (\emph{nikhārnā}) is inherited: CDIAL lists the Hindi cousin \emph{nikhārnā}, \textquotedblleft{}to clarify, to make settle\textquotedblright{}, under \emph{*niḥkṣārayati}, a form it reconstructs (CDIAL 7098). As water is left to settle until it runs clear.
\end{cousinweb}

\begin{grammarlens}[title={What you've built}]
One more word for giving reasons and opinions.
\end{grammarlens}

\subsection*{Your turn}
\begin{itemize}
\raggedright
  \item \emph{Say it:} \emph{nikhārnā}
  \item \emph{Say it:} \emph{nikhārnā}, once more
  \item \emph{Say it:} \emph{ustād ne us kā hunar nikhārā.} — The teacher brought out the best in her skill.
  \item \emph{Recall:} say the Urdu for a knack, then the Urdu for to put right, then say \emph{nikhārnā} again
\end{itemize}

\begin{tcolorbox}[breakable,colback=teal!4,colframe=teal!35!black,title={Before you move on}]
What does \ur{نکھارنا} mean? (\textquotedblleft{}To polish a skill; to bring out the best in\textquotedblright{}.) Say it once more.
\end{tcolorbox}
